\documentclass[10pt,a4paper]{amsart}
\usepackage{amsmath,amssymb,amsthm,graphicx}
% Retained minimal preamble for the Stacks Project chapter flat.tex (More on Flatness).
% Verbatim excerpt of the chapter preamble: Xy-pic setup (native preamble lines 14-18) and the native
% theorem declarations with the mathematical macros (native preamble lines 154-217). Copied mechanically;
% nothing added, removed or redefined. The wrapper class and the amsmath/amssymb/amsthm/graphicx base
% packages are supplied by the assembler, so the class/externaldocument part of the native preamble is omitted.
% For commutative diagrams we use Xy-pic
\usepackage[all]{xy}

% We use 2cell for 2-commutative diagrams.
\xyoption{2cell}
% Theorem environments.
%
\theoremstyle{plain}
\newtheorem{theorem}[subsection]{Theorem}
\newtheorem{proposition}[subsection]{Proposition}
\newtheorem{lemma}[subsection]{Lemma}

\theoremstyle{definition}
\newtheorem{definition}[subsection]{Definition}
\newtheorem{example}[subsection]{Example}
\newtheorem{exercise}[subsection]{Exercise}
\newtheorem{situation}[subsection]{Situation}

\theoremstyle{remark}
\newtheorem{remark}[subsection]{Remark}
\newtheorem{remarks}[subsection]{Remarks}

\numberwithin{equation}{subsection}

% Macros
%
\def\lim{\mathop{\mathrm{lim}}\nolimits}
\def\colim{\mathop{\mathrm{colim}}\nolimits}
\def\Spec{\mathop{\mathrm{Spec}}}
\def\Hom{\mathop{\mathrm{Hom}}\nolimits}
\def\Ext{\mathop{\mathrm{Ext}}\nolimits}
\def\SheafHom{\mathop{\mathcal{H}\!\mathit{om}}\nolimits}
\def\SheafExt{\mathop{\mathcal{E}\!\mathit{xt}}\nolimits}
\def\Sch{\mathit{Sch}}
\def\Mor{\mathop{\mathrm{Mor}}\nolimits}
\def\Ob{\mathop{\mathrm{Ob}}\nolimits}
\def\Sh{\mathop{\mathit{Sh}}\nolimits}
\def\NL{\mathop{N\!L}\nolimits}
\def\CH{\mathop{\mathrm{CH}}\nolimits}
\def\proetale{{pro\text{-}\acute{e}tale}}
\def\etale{{\acute{e}tale}}
\def\QCoh{\mathit{QCoh}}
\def\Ker{\mathop{\mathrm{Ker}}}
\def\Im{\mathop{\mathrm{Im}}}
\def\Coker{\mathop{\mathrm{Coker}}}
\def\Coim{\mathop{\mathrm{Coim}}}

% Boxtimes
%
\DeclareMathSymbol{\boxtimes}{\mathbin}{AMSa}{"02}

%
% Macros for moduli stacks/spaces
%
\def\QCohstack{\mathcal{QC}\!\mathit{oh}}
\def\Cohstack{\mathcal{C}\!\mathit{oh}}
\def\Spacesstack{\mathcal{S}\!\mathit{paces}}
\def\Quotfunctor{\mathrm{Quot}}
\def\Hilbfunctor{\mathrm{Hilb}}
\def\Curvesstack{\mathcal{C}\!\mathit{urves}}
\def\Polarizedstack{\mathcal{P}\!\mathit{olarized}}
\def\Complexesstack{\mathcal{C}\!\mathit{omplexes}}
% \Pic is the operator that assigns to X its picard group, usage \Pic(X)
% \Picardstack_{X/B} denotes the Picard stack of X over B
% \Picardfunctor_{X/B} denotes the Picard functor of X over B
\def\Pic{\mathop{\mathrm{Pic}}\nolimits}
\def\Picardstack{\mathcal{P}\!\mathit{ic}}
\def\Picardfunctor{\mathrm{Pic}}
\def\Deformationcategory{\mathcal{D}\!\mathit{ef}}

% --- appended 2026-10-12 (Chief): the project own definitions of the reference/slogan/history
% display environments, copied mechanically from the shared Stacks preamble lines 8-12
% (\usepackage{verbatim} plus the three \newenvironment lines, which the chapter preamble keeps
% verbatim from the project preamble). Nothing added, removed or redefined: these environments
% carry no mathematical content and the project itself typesets them as comments. Without these
% three lines a counted statement whose own text uses them fails with Environment slogan undefined.
\usepackage{verbatim}
\newenvironment{reference}{\comment}{\endcomment}
\newenvironment{slogan}{\comment}{\endcomment}
\newenvironment{history}{\comment}{\endcomment}
\usepackage{cleveref}
\crefname{originallemma}{original lemma}{original lemmas}
\Crefname{originallemma}{Original lemma}{Original lemmas}
\makeatletter
\newlabel{algebra-lemma-flat-universally-injective}{{058M}{source}{}{sourcepointer-1}{}}
\newlabel{algebra-lemma-flat-universally-injective@cref}{{[originallemma][0][]058M}{[1][0][]source}{}{}{}}
\makeatother
\crefname{originallemma}{original lemma}{original lemmas}
\Crefname{originallemma}{Original lemma}{Original lemmas}
\makeatletter
\newlabel{algebra-lemma-unique-prime-over-localize-below}{{00EA}{source}{}{sourcepointer-2}{}}
\newlabel{algebra-lemma-unique-prime-over-localize-below@cref}{{[originallemma][0][]00EA}{[1][0][]source}{}{}{}}
\makeatother
\crefname{originallemma}{original lemma}{original lemmas}
\Crefname{originallemma}{Original lemma}{Original lemmas}
\makeatletter
\newlabel{definition-complete-devissage-at-x-algebra}{{05HZ}{source}{}{sourcepointer-3}{}}
\newlabel{definition-complete-devissage-at-x-algebra@cref}{{[originallemma][0][]05HZ}{[1][0][]source}{}{}{}}
\makeatother
\crefname{originallemma}{original lemma}{original lemmas}
\Crefname{originallemma}{Original lemma}{Original lemmas}
\makeatletter
\newlabel{lemma-complete-devissage-flat-finite-type-module}{{05I4}{source}{}{sourcepointer-4}{}}
\newlabel{lemma-complete-devissage-flat-finite-type-module@cref}{{[originallemma][0][]05I4}{[1][0][]source}{}{}{}}
\makeatother
\crefname{originallemma}{original lemma}{original lemmas}
\Crefname{originallemma}{Original lemma}{Original lemmas}
\makeatletter
\newlabel{lemma-existence-algebra}{{05I1}{source}{}{sourcepointer-5}{}}
\newlabel{lemma-existence-algebra@cref}{{[originallemma][0][]05I1}{[1][0][]source}{}{}{}}
\makeatother
\crefname{originallemma}{original lemma}{original lemmas}
\Crefname{originallemma}{Original lemma}{Original lemmas}
\makeatletter
\newlabel{lemma-flat-at-point}{{05MN}{source}{}{sourcepointer-6}{}}
\newlabel{lemma-flat-at-point@cref}{{[originallemma][0][]05MN}{[1][0][]source}{}{}{}}
\makeatother
\crefname{originallemma}{original lemma}{original lemmas}
\Crefname{originallemma}{Original lemma}{Original lemmas}
\makeatletter
\newlabel{lemma-flat-at-point-finite}{{05P4}{source}{}{sourcepointer-7}{}}
\newlabel{lemma-flat-at-point-finite@cref}{{[originallemma][0][]05P4}{[1][0][]source}{}{}{}}
\makeatother
\crefname{originallemma}{original lemma}{original lemmas}
\Crefname{originallemma}{Original lemma}{Original lemmas}
\makeatletter
\newlabel{lemma-flat-at-point-go-up}{{05P5}{source}{}{sourcepointer-8}{}}
\newlabel{lemma-flat-at-point-go-up@cref}{{[originallemma][0][]05P5}{[1][0][]source}{}{}{}}
\makeatother
\crefname{originallemma}{original lemma}{original lemmas}
\Crefname{originallemma}{Original lemma}{Original lemmas}
\makeatletter
\newlabel{lemma-induction-step}{{05I3}{source}{}{sourcepointer-9}{}}
\newlabel{lemma-induction-step@cref}{{[originallemma][0][]05I3}{[1][0][]source}{}{}{}}
\makeatother
\crefname{originallemma}{original lemma}{original lemmas}
\Crefname{originallemma}{Original lemma}{Original lemmas}
\makeatletter
\newlabel{more-morphisms-lemma-finite-morphism-single-point-in-fibre}{{05B8}{source}{}{sourcepointer-10}{}}
\newlabel{more-morphisms-lemma-finite-morphism-single-point-in-fibre@cref}{{[originallemma][0][]05B8}{[1][0][]source}{}{}{}}
\makeatother
\crefname{originallemma}{original lemma}{original lemmas}
\Crefname{originallemma}{Original lemma}{Original lemmas}
\makeatletter
\newlabel{situation-flat-at-point}{{05MK}{source}{}{sourcepointer-11}{}}
\newlabel{situation-flat-at-point@cref}{{[originallemma][0][]05MK}{[1][0][]source}{}{}{}}
\makeatother
\crefname{originallemma}{original lemma}{original lemmas}
\Crefname{originallemma}{Original lemma}{Original lemmas}
\makeatletter
\newlabel{theorem-flattening-map}{{05PF}{source}{}{sourcepointer-12}{}}
\newlabel{theorem-flattening-map@cref}{{[originallemma][0][]05PF}{[1][0][]source}{}{}{}}
\makeatother
\title{Corepresentability of the flat at a point functor in the henselian essentially finite type case}
\author{The Stacks Project Authors}
\date{stacks-project:flat.tex:local-snapshot-20261011; GNU FDL 1.2 or later}
\begin{document}
\maketitle
\noindent\textit{Editorial note.} All delivered mathematical text is the author's own, copied from the pinned Stacks Project chapter file flat.tex; no statement or proof step was rewritten, reordered or removed.\par\medskip
\noindent\textit{Editorial note.} Preamble provenance: the chapter's own retained minimal preamble (sources/stacks-flat/retained-preamble.tex, lines 1-85) is a verbatim excerpt of the native Xy-pic setup and of the native theorem declarations and mathematical macros, followed by the project's own reference/slogan/history display-environment definitions copied mechanically from the shared preamble; the wrapper class and the amsmath/amssymb/amsthm/graphicx base packages are supplied by the assembler, and extra\_packages is left empty because everything needed is already in the retained preamble. No macro was invented or redefined.\par\medskip
\noindent\textit{Editorial note.} Same-chapter pointers are located at their \textbackslash{}label lines inside the pinned chapter file; cross-chapter pointers are located in the project's own tag list staged in the snapshot (sources/stacks-chow/tags.txt), at the verbatim lines carrying their canonical tags.\par\medskip
\noindent\textit{Allowed standard prerequisites (editorial).}\begin{itemize}
\item The functor of points where a module is flat, and its invariance under a flat local extension of the target. Same chapter flat.tex, situation labelled situation-flat-at-point (statement, proof line 6792), labels lemma-flat-at-point-go-up (proof line 6816), definition-complete-devissage-at-x-algebra (proof line 6828), lemma-flat-at-point-finite (proof lines 6831 and 6840) and lemma-flat-at-point part (2) (proof line 6900). Fixes the functor F\_lf being corepresented, makes it unchanged when B is replaced by a local ring flat over B, transports the problem along the finite and surjective steps of a complete devissage, and converts the corepresentation for a filtered union of ideals into finite generation.
\item Existence of a complete devissage in the henselian case. Same chapter flat.tex, label lemma-existence-algebra, invoked at proof line 6819. Allows one to assume a complete devissage of N over C over A at q; this is the structural input that makes the induction on the length of the devissage possible.
\item Descent of flatness along the induction step and flatness of a universally exact sequence. Same chapter flat.tex, labels lemma-induction-step (proof lines 6853 and 6874), lemma-complete-devissage-flat-finite-type-module (proof line 6871) and theorem-flattening-map (proof line 6889); Algebra, label algebra-lemma-flat-universally-injective (proof line 6885); Algebra, label algebra-lemma-unique-prime-over-localize-below and More on Morphisms, label more-morphisms-lemma-finite-morphism-single-point-in-fibre (proof lines 6836-6838). Gives that flatness propagates from M\_1 to its cokernel along the induction step, makes the cokernel free over A once it is flat, turns the exact sequence into a universally exact one so that injectivity of the base changed map is detected on the presentation, and supplies the flattering criterion of the preceding theorem for the final corepresentation.
\end{itemize}\medskip
\ifdefined\hypertarget\hypertarget{sourcepointer-1}{}\fi
\noindent\textit{Source reference.} Original lemma 058M: sources/stacks-chow/tags.txt, source lines 6785--6785. Navigation pointer for \textbackslash{}ref\{algebra-lemma-flat-universally-injective\} used inside the counted span (line 6885); the argument that consumes it is retained in the span.\par\medskip
\ifdefined\hypertarget\hypertarget{sourcepointer-2}{}\fi
\noindent\textit{Source reference.} Original lemma 00EA: sources/stacks-chow/tags.txt, source lines 527--527. Navigation pointer for \textbackslash{}ref\{algebra-lemma-unique-prime-over-localize-below\} used inside the counted span (line 6836); the argument that consumes it is retained in the span.\par\medskip
\ifdefined\hypertarget\hypertarget{sourcepointer-3}{}\fi
\noindent\textit{Source reference.} Original lemma 05HZ: sources/stacks-flat/pinned-original-source.tex, source lines 1710--1710. Navigation pointer for \textbackslash{}ref\{definition-complete-devissage-at-x-algebra\} used inside the counted span (line 6828); the argument that consumes it is retained in the span.\par\medskip
\ifdefined\hypertarget\hypertarget{sourcepointer-4}{}\fi
\noindent\textit{Source reference.} Original lemma 05I4: sources/stacks-flat/pinned-original-source.tex, source lines 2426--2426. Navigation pointer for \textbackslash{}ref\{lemma-complete-devissage-flat-finite-type-module\} used inside the counted span (line 6871); the argument that consumes it is retained in the span.\par\medskip
\ifdefined\hypertarget\hypertarget{sourcepointer-5}{}\fi
\noindent\textit{Source reference.} Original lemma 05I1: sources/stacks-flat/pinned-original-source.tex, source lines 1757--1757. Navigation pointer for \textbackslash{}ref\{lemma-existence-algebra\} used inside the counted span (line 6819); the argument that consumes it is retained in the span.\par\medskip
\ifdefined\hypertarget\hypertarget{sourcepointer-6}{}\fi
\noindent\textit{Source reference.} Original lemma 05MN: sources/stacks-flat/pinned-original-source.tex, source lines 6045--6045. Navigation pointer for \textbackslash{}ref\{lemma-flat-at-point\} used inside the counted span (line 6900); the argument that consumes it is retained in the span.\par\medskip
\ifdefined\hypertarget\hypertarget{sourcepointer-7}{}\fi
\noindent\textit{Source reference.} Original lemma 05P4: sources/stacks-flat/pinned-original-source.tex, source lines 6068--6068. Navigation pointer for \textbackslash{}ref\{lemma-flat-at-point-finite\} used inside the counted span (line 6831); the argument that consumes it is retained in the span.\par\medskip
\ifdefined\hypertarget\hypertarget{sourcepointer-8}{}\fi
\noindent\textit{Source reference.} Original lemma 05P5: sources/stacks-flat/pinned-original-source.tex, source lines 6108--6108. Navigation pointer for \textbackslash{}ref\{lemma-flat-at-point-go-up\} used inside the counted span (line 6816); the argument that consumes it is retained in the span.\par\medskip
\ifdefined\hypertarget\hypertarget{sourcepointer-9}{}\fi
\noindent\textit{Source reference.} Original lemma 05I3: sources/stacks-flat/pinned-original-source.tex, source lines 2356--2356. Navigation pointer for \textbackslash{}ref\{lemma-induction-step\} used inside the counted span (line 6853); the argument that consumes it is retained in the span.\par\medskip
\ifdefined\hypertarget\hypertarget{sourcepointer-10}{}\fi
\noindent\textit{Source reference.} Original lemma 05B8: sources/stacks-chow/tags.txt, source lines 6876--6876. Navigation pointer for \textbackslash{}ref\{more-morphisms-lemma-finite-morphism-single-point-in-fibre\} used inside the counted span (line 6838); the argument that consumes it is retained in the span.\par\medskip
\ifdefined\hypertarget\hypertarget{sourcepointer-11}{}\fi
\noindent\textit{Source reference.} Original lemma 05MK: sources/stacks-flat/pinned-original-source.tex, source lines 6008--6008. Navigation pointer for \textbackslash{}ref\{situation-flat-at-point\} used inside the counted span (line 6792); the argument that consumes it is retained in the span.\par\medskip
\ifdefined\hypertarget\hypertarget{sourcepointer-12}{}\fi
\noindent\textit{Source reference.} Original lemma 05PF: sources/stacks-flat/pinned-original-source.tex, source lines 6671--6671. Navigation pointer for \textbackslash{}ref\{theorem-flattening-map\} used inside the counted span (line 6889); the argument that consumes it is retained in the span.\par\medskip
\begin{theorem}
\label{theorem-flattening-local}
In
Situation \ref{situation-flat-at-point}
assume $A$ is henselian, $B$ is essentially of finite type over $A$, and
$M$ is a finite $B$-module. Then there exists an ideal
$I \subset A$ such that $A/I$ corepresents the functor $F_{lf}$ on the category
$\mathcal{C}$. In other words given a local homomorphism of local rings
$\varphi : A \to A'$ with $B' = B \otimes_A A'$ and $M' = M \otimes_A A'$
the following are equivalent:
\begin{enumerate}
\item $\forall \mathfrak q \in V(\mathfrak m_{A'}B' + \mathfrak m_B B')
\subset \Spec(B') :
M'_{\mathfrak q}\text{ is flat over }A'$, and
\item $\varphi(I) = 0$.
\end{enumerate}
If $B$ is essentially of finite presentation over $A$ and $M$
of finite presentation over $B$, then $I$ is a finitely generated ideal.
\end{theorem}

\begin{proof}
Choose a finite type ring map $A \to C$ and a finite $C$-module $N$
and a prime $\mathfrak q$ of $C$ such that $B = C_{\mathfrak q}$
and $M = N_{\mathfrak q}$. In the following, when we say
``the theorem holds for $(N/C/A, \mathfrak q)$ we mean that
it holds for $(A \to B, M)$ where $B = C_{\mathfrak q}$ and
$M = N_{\mathfrak q}$. By
Lemma \ref{lemma-flat-at-point-go-up}
the functor $F_{lf}$ is unchanged if we replace $B$ by a local ring
flat over $B$. Hence, since $A$ is henselian, we may apply
Lemma \ref{lemma-existence-algebra}
and assume that there exists a complete d\'evissage of
$N/C/A$ at $\mathfrak q$.

\medskip\noindent
Let $(A_i, B_i, M_i, \alpha_i, \mathfrak q_i)_{i = 1, \ldots, n}$
be such a complete d\'evissage of $N/C/A$ at $\mathfrak q$. Let
$\mathfrak q'_i \subset A_i$ be the unique prime lying over
$\mathfrak q_i \subset B_i$ as in
Definition \ref{definition-complete-devissage-at-x-algebra}.
Since $C \to A_1$ is surjective and $N \cong M_1$ as $C$-modules,
we see by
Lemma \ref{lemma-flat-at-point-finite}
it suffices to prove the theorem holds for $(M_1/A_1/A, \mathfrak q'_1)$.
Since $B_1 \to A_1$ is finite and $\mathfrak q_1$ is the only prime
of $B_1$ over $\mathfrak q'_1$ we see that
$(A_1)_{\mathfrak q'_1} \to (B_1)_{\mathfrak q_1}$ is finite (see
Algebra, Lemma \ref{algebra-lemma-unique-prime-over-localize-below} or
More on Morphisms,
Lemma \ref{more-morphisms-lemma-finite-morphism-single-point-in-fibre}).
Hence by
Lemma \ref{lemma-flat-at-point-finite}
it suffices to prove the theorem holds for $(M_1/B_1/A, \mathfrak q_1)$.

\medskip\noindent
At this point we may assume, by induction on the length $n$ of the
d\'evissage, that the theorem holds for $(M_2/B_2/A, \mathfrak q_2)$.
(If $n = 1$, then $M_2 = 0$ which is flat over $A$.)
Reversing the last couple of steps of the previous paragraph, using
that $M_2 \cong \Coker(\alpha_2)$ as $B_1$-modules, we see
that the theorem holds for $(\Coker(\alpha_1)/B_1/A, \mathfrak q_1)$.

\medskip\noindent
Let $A'$ be an object of $\mathcal{C}$. At this point we use
Lemma \ref{lemma-induction-step}
to see that if $(M_1 \otimes_A A')_{\mathfrak q'}$ is flat
over $A'$ for a prime $\mathfrak q'$ of $B_1 \otimes_A A'$
lying over $\mathfrak m_{A'}$, then
$(\Coker(\alpha_1) \otimes_A A')_{\mathfrak q'}$ is flat over $A'$.
Hence we conclude that $F_{lf}$ is a subfunctor of the
functor $F'_{lf}$ associated to the module
$\Coker(\alpha_1)_{\mathfrak q_1}$ over $(B_1)_{\mathfrak q_1}$.
By the previous paragraph we know $F'_{lf}$ is corepresented by
$A/J$ for some ideal $J \subset A$. Hence we may replace $A$ by
$A/J$ and assume that $\Coker(\alpha_1)_{\mathfrak q_1}$ is
flat over $A$.

\medskip\noindent
Since $\Coker(\alpha_1)$ is a $B_1$-module for which
there exist a complete d\'evissage of $N_1/B_1/A$ at $\mathfrak q_1$
and since $\Coker(\alpha_1)_{\mathfrak q_1}$ is
flat over $A$ by
Lemma \ref{lemma-complete-devissage-flat-finite-type-module}
we see that $\Coker(\alpha_1)$ is free as an $A$-module, in particular
flat as an $A$-module. Hence
Lemma \ref{lemma-induction-step}
implies $F_{lf}(A')$ is nonempty if and only if $\alpha \otimes 1_{A'}$
is injective. Let $N_1 = \Im(\alpha_1) \subset M_1$ so that
we have exact sequences
$$
0 \to N_1 \to M_1 \to \Coker(\alpha_1) \to 0
\quad\text{and}\quad
B_1^{\oplus r_1} \to N_1 \to 0
$$
The flatness of $\Coker(\alpha_1)$ implies the first sequence
is universally exact (see
Algebra, Lemma \ref{algebra-lemma-flat-universally-injective}).
Hence $\alpha \otimes 1_{A'}$ is injective if and only if
$B_1^{\oplus r_1} \otimes_A A' \to N_1 \otimes_A A'$
is an isomorphism. Finally,
Theorem \ref{theorem-flattening-map}
applies to show this functor is corepresentable by $A/I$ for some ideal $I$
and we conclude $F_{lf}$ is corepresentable by $A/I$ also.

\medskip\noindent
To prove the final statement, suppose that $A \to B$ is essentially of finite
presentation and $M$ of finite presentation over $B$. Let $I \subset A$
be the ideal such that $F_{lf}$ is corepresented by $A/I$.
Write $I = \bigcup I_\lambda$ where $I_\lambda$ ranges over the finitely
generated ideals contained in $I$. Then, since $F_{lf}(A/I) = \{*\}$
we see that $F_{lf}(A/I_\lambda) = \{*\}$ for some $\lambda$, see
Lemma \ref{lemma-flat-at-point} part (2).
Clearly this implies that $I = I_\lambda$.
\end{proof}

\end{document}
